% -----------------------------------------------------------------------------
% R3: degenerate stationary-hull residual class
% -----------------------------------------------------------------------------
% This section is written as a rigorous PDE insert.  It is intentionally
% conservative.  The naive formulation "the time-translation hull merely meets a
% bifurcating stationary family" is too broad: an inactive stationary profile in
% the hull does not by itself contradict singularity formation.  The useful
% residual class is the active-attained degenerate stationary-hull class: the
% hull must contain a stationary-family profile which is itself an active
% Seregin limit, i.e. it retains a positive local CKN concentration.
%
% The section proves the exact implication needed to close this class:
%
%   no active Seregin-attainable degenerate stationary-family profile
%       ==> R_3 = empty.
%
% It also proves the standard structural facts used in that reduction: hull
% limits of Type I Seregin profiles are again Type I Seregin profiles after
% reselecting scales; nontrivial C^1 stationary families produce linearized
% kernel elements; and any exact trajectory constrained to a stationary family
% is stationary.
% -----------------------------------------------------------------------------

\section{The degenerate stationary-hull residual class}
\label{sec:R3-degenerate-stationary-hull}

We now isolate the third refined residual class.  This class is designed to
capture Type~I ancient limits whose time-translation hull contains an active
stationary profile lying on a non-isolated family of stationary centered
profiles.  The role of the stationary-family condition is to separate genuinely degenerate
stationary profile families from isolated stationary profiles and from the
symmetry-generated classes already removed in \(\Cax\) and \(\Crot\).

A small point of formulation is important.  It is not enough to assume that the
hull of a residual ancient solution contains some stationary-family element.  A
stationary element may occur as an inactive limit and need not retain the local
CKN concentration inherited from the original singular point.  Such an inactive
hull element should not be placed in a class that is to be closed by a
stationary-profile obstruction.  Therefore the class below is defined using
\emph{active attainment}: the stationary-family profile must itself occur as a
Seregin limit carrying a positive local CKN concentration.

Throughout this section the centered equation is
\begin{equation}\label{eq:centered-R3}
   \partial_\tau V-\Delta V+\frac12 y\cdot\nabla V+\frac12 V
   +(V\cdot\nabla)V+\nabla P=0,
   \qquad \nabla\cdot V=0 .
\end{equation}
Its projected form is
\begin{equation}\label{eq:projected-centered-R3}
   \partial_\tau V=\mathcal F(V),
   \qquad
   \mathcal F(W)
   :=
   \mathbb P\left[
      \Delta W-\frac12(W+y\cdot\nabla W)-(W\cdot\nabla)W
   \right],
\end{equation}
where \(\mathbb P\) is the whole-space Leray projector.

\subsection{Stationary branches and active hull attainment}

\begin{definition}[Admissible stationary phase space]
\label{def:R3-stationary-phase-space}
An admissible stationary phase space is a Banach space \(\mathfrak B\) of
solenoidal vector fields on \(\mathbb R^3\) such that:
\begin{enumerate}[label=(\roman*)]
   \item \(\mathfrak B\hookrightarrow C^{2,\alpha}_{\loc}(\mathbb R^3)
      \cap L^\infty(\mathbb R^3)\) continuously for some
      \(\alpha\in(0,1)\);
   \item the map \(\mathcal F\) in \eqref{eq:projected-centered-R3} is
      well-defined on an open subset \(\mathcal U\subset\mathfrak B\) and maps
      \(\mathcal U\) into a Banach space \(\mathfrak Y\);
   \item \(\mathcal F:\mathcal U\to\mathfrak Y\) is \(C^1\).
\end{enumerate}
\end{definition}

\begin{definition}[Nontrivial stationary family]
\label{def:R3-stationary-family}
Let \(W\in\mathcal U\subset\mathfrak B\) satisfy \(\mathcal F(W)=0\).  We say
that \(W\) lies on an admissible nontrivial stationary family if there exist
\(\varepsilon>0\) and a \(C^1\) curve
\[
   a\mapsto W_a\in\mathcal U,
   \qquad |a|<\varepsilon,
\]
such that
\begin{equation}\label{eq:R3-stationary-family}
   W_0=W,
   \qquad
   \mathcal F(W_a)=0 \quad (|a|<\varepsilon),
\end{equation}
and such that the family is not locally generated only by the action of the
rigid rotation group.  Equivalently, after possibly reparametrizing the curve,
one may assume that
\begin{equation}\label{eq:R3-nontrivial-tangent}
   \dot W_0\notin
   \{\Omega W-(\Omega y)\cdot\nabla W:\Omega\in\mathfrak{so}(3)\}
\end{equation}
whenever the tangent \(\dot W_0\) exists and is nonzero.
\end{definition}

\begin{definition}[Local CKN activity]
\label{def:R3-active-profile}
Let \((W,\Pi)\) be a stationary centered profile.  We say that \((W,\Pi)\) is
\(\varepsilon_*\)-active if there exist a compact parabolic cylinder
\(Q_r(z_0)\) and an admissible pressure gauge \(c_{r,z_0}\) such that
\begin{equation}\label{eq:R3-active-ckn}
   \iint_{Q_r(z_0)}
   \left(|W|^3+|\Pi-c_{r,z_0}|^{3/2}\right)\,dy\,d\tau
   \ge \varepsilon_* .
\end{equation}
For a stationary profile the integrand is independent of \(\tau\), but the
parabolic-cylinder notation is retained to match the CKN normalization used for
Seregin limits.
\end{definition}

\begin{definition}[Active degenerate stationary-hull residual class]
\label{def:R3-active-class}
Fix an admissible stationary phase space \(\mathfrak B\).  The active
degenerate stationary-hull residual class
\[
   \Cstat(\mathcal S;I,J;\mathfrak B)
\]
consists of those \(V\in\mathcal R(\mathcal S;I,J)\), after the ordered
removal of \(\Cax\) and \(\Crot\), for which there exists a
sequence \(s_k\in\mathbb R\) and a stationary profile \((W,\Pi)\) such that
\begin{equation}\label{eq:R3-hull-convergence}
   (V(\cdot,\cdot+s_k),P(\cdot,\cdot+s_k)-c_k(\cdot))
   \longrightarrow (W,\Pi)
\end{equation}
locally smoothly for the velocity and locally in \(L^{3/2}\) for the pressure
modulo admissible gauges, and such that:
\begin{enumerate}[label=(\roman*)]
   \item \(W\in\mathfrak B\) and \(\mathcal F(W)=0\);
   \item \(W\) lies on an admissible nontrivial stationary family in the sense
      of \Cref{def:R3-stationary-family};
   \item \((W,\Pi)\) is \(\varepsilon_*\)-active in the sense of
      \Cref{def:R3-active-profile};
   \item \((W,\Pi)\) is not in any previously closed stationary, structured,
      axisymmetric bounded-circulation, or rotational relative-equilibrium
      class.
\end{enumerate}
In the ordered residual decomposition one keeps this active-attainment class as
the degenerate stationary-hull alternative.  The inactive hull cases
should be assigned to the generic terminal residual class, where they are
handled by the terminal local concentration decomposition rather than by a
stationary-profile obstruction.
\end{definition}

\begin{remark}[Why the active condition is necessary]
\label{rem:R3-active-necessary}
The condition \eqref{eq:R3-active-ckn} prevents the following logical error.  A
nonstationary ancient solution may have an inactive stationary profile in its
hull, for instance a profile with zero local CKN mass on the relevant compact
cylinders.  Excluding the stationary profile would not exclude the original
ancient solution.  The stationary-profile obstruction used below applies only to stationary
profiles that are themselves attainable as nontrivial Seregin limits with
positive local concentration.
\end{remark}

\subsection{Hull limits are again Seregin limits}

The next lemma is the basic compatibility between time-translation hulls in
centered variables and rescaling sequences in physical variables.  It is the
reason that an active stationary profile in the hull may be treated as an
attained blow-up profile, not merely as an internal dynamical limit of the
already extracted ancient solution.

\begin{lemma}[Scale reselection for centered hull limits]
\label{lem:R3-hull-is-seregin}
Let \((V,P)\) be a centered Type~I Seregin limit obtained from a finite-energy
suitable weak solution \((u,p)\) at a singular point \(z_*=(x_*,t_*)\).  Suppose
that, for some sequence \(s_k\in\mathbb R\),
\[
   (V(\cdot,\cdot+s_k),P(\cdot,\cdot+s_k)-c_k(\cdot))
   \to (W,\Pi)
\]
locally in the Seregin topology.  Then, after passing to a diagonal subsequence
and reselecting the physical blow-up scales, \((W,\Pi)\) is itself a centered
Type~I Seregin limit of \((u,p)\) at \(z_*\).  If the convergence to \((W,\Pi)\)
retains a positive local CKN lower bound, then \((W,\Pi)\) is a nontrivial
Seregin limit with the same type of inherited nonvanishing condition.
\end{lemma}

\begin{proof}
Let \(\lambda_n\downarrow0\) be a sequence of scales producing \((V,P)\).  In
centered variables, a time translation by \(s\) is exactly a change of the
physical blow-up scale by a factor depending only on \(s\).  With the convention
used in \eqref{eq:centered-R3}, write this adjusted scale as
\begin{equation}\label{eq:R3-adjusted-scale}
   \lambda_{n,s}:=e^{-s/2}\lambda_n .
\end{equation}
The rescaled fields built from \((u,p)\) at scale \(\lambda_{n,s}\) are the same
as the fields at scale \(\lambda_n\), translated by \(s\) in the centered time
variable, up to the corresponding pressure gauge.

Fix an exhaustion of compact cylinders \(K_m\subset\mathbb R^3\times\mathbb R\).
For each \(k\), choose \(n(k)\) so large that \(\lambda_{n(k),s_k}\to0\) and the
rescaled fields at scale \(\lambda_{n(k)}\), shifted by \(s_k\), are within
\(1/k\) of \((V(\cdot,\cdot+s_k),P(\cdot,\cdot+s_k)-c_k)\) on \(K_k\) in the
local Seregin topology.  This is possible because \(\lambda_n\downarrow0\) and
because the original sequence converges locally to \((V,P)\).

By the assumed hull convergence, the shifted limit
\((V(\cdot,\cdot+s_k),P(\cdot,\cdot+s_k)-c_k)\) converges to \((W,\Pi)\) on each
fixed compact cylinder.  The diagonal sequence of physical rescalings with
scales \(\lambda_{n(k),s_k}\) therefore converges locally to \((W,\Pi)\).  Hence
\((W,\Pi)\) is again a centered Type~I Seregin limit at \(z_*\).  The assertion
about nonvanishing follows from lower semicontinuity and local convergence of
the CKN quantities after the prescribed pressure gauges.
\end{proof}

\subsection{Elementary consequences of stationary-family degeneracy}

The stationary-family condition has two immediate PDE consequences.  First, a differentiable
nontrivial family produces a nontrivial element of the kernel of the linearized
stationary Leray operator.  Second, if an ancient trajectory is exactly
constrained to a stationary family, then it is stationary.

\begin{lemma}[Stationary family implies linearized degeneracy]
\label{lem:R3-family-kernel}
Let \(a\mapsto W_a\in\mathfrak B\) be a \(C^1\) stationary family satisfying
\eqref{eq:R3-stationary-family}.  Then
\begin{equation}\label{eq:R3-kernel}
   D\mathcal F(W_0)\dot W_0=0 .
\end{equation}
In particular, if \(\dot W_0\ne0\) modulo infinitesimal rotations, the
linearized stationary Leray operator has a genuine non-symmetry kernel element.
\end{lemma}

\begin{proof}
Differentiate \(\mathcal F(W_a)=0\) at \(a=0\).  Since \(\mathcal F\) is \(C^1\)
on \(\mathfrak B\), this gives \eqref{eq:R3-kernel}.  The final statement is
exactly the nontriviality condition in \Cref{def:R3-stationary-family}.
\end{proof}

\begin{lemma}[Trajectories constrained to a stationary family are stationary]
\label{lem:R3-family-constrained-stationary}
Let \(a\mapsto W_a\in\mathfrak B\) be a \(C^1\) stationary family with
\(\dot W_a\ne0\) on an interval.  Suppose that a smooth solution of
\eqref{eq:projected-centered-R3} has the form
\begin{equation}\label{eq:R3-family-trajectory}
   V(\tau)=W_{a(\tau)}
\end{equation}
for a \(C^1\) function \(a(\tau)\).  Then \(a\) is constant, and hence
\(V\) is stationary.
\end{lemma}

\begin{proof}
Since \(\mathcal F(W_a)=0\) for every \(a\), equation
\eqref{eq:projected-centered-R3} gives
\[
   a'(\tau)\dot W_{a(\tau)}=\partial_\tau V(\tau)=\mathcal F(W_{a(\tau)})=0 .
\]
Because \(\dot W_{a(\tau)}\ne0\), we obtain \(a'(\tau)=0\).  Thus \(a\) is
constant on the interval.
\end{proof}

\begin{corollary}[Exact stationary-family subcase]
\label{cor:R3-exact-family-subcase}
If an element of the residual class is itself an exact trajectory on an
admissible stationary family, then it is a stationary centered profile.  Hence
it is excluded whenever it satisfies the stationary \(L^3\) Liouville
criterion, or more generally whenever the stationary profile belongs to one of
the previously closed stationary or structured classes.
\end{corollary}

\begin{proof}
Let \(V(\tau)=W_{a(\tau)}\) be such an exact stationary-family trajectory.  By
\Cref{lem:R3-family-constrained-stationary}, \(a'(\tau)=0\) on every interval on
which the stationary-family parametrization is nondegenerate.  Hence \(a\) is constant and
\(V\) is a single stationary centered profile.  If that stationary profile
satisfies the stationary \(L^3\) Liouville criterion or lies in one of the
previously closed structured or ordered lower residual classes, the corresponding
prior exclusion removes it from the residual class.
\end{proof}

\subsection{The exact attainability statement for the stationary-hull class}

The remaining issue is not the internal dynamics along a stationary family; exact
stationary-family
motion is stationary by \Cref{lem:R3-family-constrained-stationary}.  The real
issue is attainability: can a degenerate stationary-family profile occur as an
active Seregin limit of a finite-energy Type~I source?  The following statement
is the precise attainability statement required to close the active
degenerate stationary-hull class.  It is proved in the following section by
the terminal local concentration theorem.

\begin{theorem}[No active degenerate stationary-family profile is Seregin-attainable]
\label{thm:R3-no-attainable-degenerate-family}
Fix an admissible stationary phase space \(\mathfrak B\).  Let \((W,\Pi)\) be a
stationary centered profile satisfying:
\begin{enumerate}[label=(\roman*)]
   \item \(W\in\mathfrak B\), \(\mathcal F(W)=0\), and \(W\) lies on an
      admissible nontrivial stationary family;
   \item \((W,\Pi)\) is a centered Type~I Seregin limit of a finite-energy
      suitable weak solution at a singular point;
   \item \((W,\Pi)\) is \(\varepsilon_*\)-active;
   \item \((W,\Pi)\) is not in the small, stationary \(L^3\), tight,
      fast-decay, structured, axisymmetric bounded-circulation, or rotational
      relative-equilibrium classes already removed earlier in the stratification.
\end{enumerate}
Then no such \((W,\Pi)\) exists.
\end{theorem}

\begin{proof}
This is exactly \Cref{cor:R3S-no-active-degenerate-family}, proved below from
the local terminal stratification theorem.
\end{proof}

\begin{remark}[Interpretation of the attainability theorem]
\label{rem:R3-input-interpretation}
\Cref{thm:R3-no-attainable-degenerate-family} is a forward-attainability
obstruction.  It does not assert that degenerate stationary profiles cannot
exist as elliptic solutions of the stationary centered equation.  It asserts
only that such profiles cannot occur as active Type~I blow-up limits of
finite-energy suitable weak solutions after all previously closed classes have
been removed.  This is the correct formulation for the residual problem: the
elliptic family may exist, but it must be inaccessible to the local blow-up
procedure coming from a finite-energy Leray source.
\end{remark}

\subsection{Closure of the active degenerate stationary-hull class}

\begin{theorem}[Exclusion of the active degenerate stationary-hull class]
\label{thm:R3-exclusion}
\begin{equation}\label{eq:R3-empty}
   \Cstat(\mathcal S;I,J;\mathfrak B)=\varnothing .
\end{equation}
Consequently the ordered third residual class is empty, provided it is defined
with the active-attainment condition in \Cref{def:R3-active-class}.
\end{theorem}

\begin{proof}
Suppose, for contradiction, that
\(V\in\Cstat(\mathcal S;I,J;\mathfrak B)\).  By
\Cref{def:R3-active-class}, there is a stationary profile \((W,\Pi)\) in the
centered time-translation hull of \((V,P)\), and \((W,\Pi)\) is
\(\varepsilon_*\)-active.  The profile \(W\) belongs to \(\mathfrak B\), satisfies
\(\mathcal F(W)=0\), and lies on an admissible nontrivial stationary family.
Moreover, by the ordered definition of \(\Cstat\), it is not
in any of the previously closed lower classes.

By \Cref{lem:R3-hull-is-seregin}, the hull profile \((W,\Pi)\) is itself a
centered Type~I Seregin limit of the original finite-energy suitable weak
solution at the same singular point, after reselecting scales.  Since activity
is retained, \((W,\Pi)\) satisfies every condition in
\Cref{thm:R3-no-attainable-degenerate-family}.  This contradicts the
attainability theorem.
Therefore no such \(V\) exists, proving \eqref{eq:R3-empty}.
\end{proof}

\begin{proposition}[Already-closed integrable stationary-family subcase]
\label{prop:R3-L3-subcase}
Let \((W,\Pi)\) be a stationary centered profile lying on an admissible
stationary family.  If
\begin{equation}\label{eq:R3-L3}
   W\in L^3(\mathbb R^3),
\end{equation}
then \(W\equiv0\).  Hence no active hull profile satisfying
\eqref{eq:R3-L3} can occur in \(\Cstat\).
\end{proposition}

\begin{proof}
The stationary physical representative associated with \(W\) is a backward
self-similar Navier--Stokes solution.  The assumption \(W\in L^3(\mathbb R^3)\)
places it in the stationary \(L^3\) class already excluded by the classical
Ne\v cas--R\r u\v zi\v cka--\v Sver\'ak Liouville theorem, equivalently by the
endpoint ancient Liouville theorem cited in Paper I.  Hence \(W\equiv0\).  A
zero profile has zero pressure-normalized local CKN mass after choosing the
admissible pressure gauge, so it cannot be \(\varepsilon_*\)-active.
\end{proof}

\begin{remark}[Closure after terminal local concentration analysis]
\label{rem:R3-final-obligation}
The preceding argument reduces the active stationary-hull class from a dynamical hull question to a static
attainability question for stationary centered profiles in nontrivial stationary
families.
\Cref{thm:R3-no-attainable-degenerate-family} is supplied by the terminal local
concentration analysis in the next section.
\end{remark}

\begin{remark}[Recommended statement for the residual decomposition]
\label{rem:R3-decomposition-replacement}
For the residual decomposition, the third closure line should be stated as follows:
\[
   \text{active degenerate stationary-hull residual}
   \quad\xrightarrow{\;\text{no active stationary-family attainability}\;}\quad
   \varnothing .
\]
The older phrase ``bifurcating profile families'' should be used only as
informal motivation.  The theorem-level class should be called the
\emph{active degenerate stationary-hull residual class}.
\end{remark}
